% =====================================================================
% abstract.tex -- content of the MCM/ICM Summary Sheet (page 1).
%
% It defines two commands used by \maketitle:
%   \summary{...}   one or two paragraphs: background -> methods -> key
%                   numeric results for each sub-problem -> conclusion
%   \keywords{...}  3-6 keywords, comma separated
%
% Ownership: normally maintained by the part/hNN/misc branch, because the
% summary integrates all problems. It must be written/updated LAST, after
% p1/p2/p3/sens are frozen.
%
% Style rules: no formulas, no citations, no undefined abbreviations, and
% every claim must be backed by a number that also appears in the body.
% =====================================================================

\summary{%
Write one to two paragraphs here. Start with the background and the task.
Then, for each sub-problem, state the model or method used and the key
result (with concrete numbers). Close with the main conclusion and, if
relevant, one sentence on sensitivity/robustness. Replace every sentence
of this placeholder before submitting; the judges read this page first.

A second paragraph can describe the validation of the model and the
practical implication of the results.
}

\keywords{keyword one, keyword two, keyword three, keyword four}
